%=====================================================================
\chapter{High Availability, Live Migration and Disaster Recovery}\label{ch:migration}
%=====================================================================
\locator{5}{Part V --- Management, Performance and Security}

\begin{outcomes}
By the end of this chapter you will be able to:
\begin{tight}
  \item \textbf{Describe} pre-copy live migration round by round, and
        \textbf{state} the condition under which it converges.
  \item \textbf{Compute} the number of rounds and the total transfer for a
        given guest, dirty rate and link speed.
  \item \textbf{Explain} what happens when it does not converge, and
        \textbf{compare} auto-convergence with post-copy.
  \item \textbf{Distinguish} high availability from fault tolerance from
        disaster recovery, and \textbf{express} each as RPO and RTO.
  \item \textbf{Compute} the downtime budget implied by an availability target,
        and \textbf{say} what each additional nine costs.
\end{tight}
\end{outcomes}

\begin{prereq}
\chref{cpumem} (memory, and the fact that a running guest's pages change under
you) and \chref{storage} (shared storage, without which almost nothing here
works). \chref{paravirt}'s warning that an assigned device cannot be migrated
is the exception that proves this chapter's rule.
\end{prereq}

\begin{keyterms}
live migration $\bullet$ pre-copy $\bullet$ post-copy $\bullet$ dirty page
$\bullet$ dirty rate $\bullet$ convergence $\bullet$ switchover $\bullet$
stop-and-copy $\bullet$ auto-convergence $\bullet$ page fault over the network
$\bullet$ high availability $\bullet$ fault tolerance $\bullet$ cluster
$\bullet$ quorum $\bullet$ split brain $\bullet$ fencing $\bullet$ STONITH
$\bullet$ RPO $\bullet$ RTO $\bullet$ availability $\bullet$ nines $\bullet$
3-2-1 rule $\bullet$ immutable backup
\end{keyterms}

\section{Moving a Running Machine}

Nothing else in computing does this. A guest serving traffic moves from one
physical host to another, keeps its memory, its open connections and its
in-flight work, and pauses for a time its users cannot perceive. It is the
capability that makes everything in \chref{management} possible: hosts can be
patched, rebalanced and retired without anyone scheduling an outage.

The difficulty is stated in one sentence: \emph{the memory you are copying is
being modified while you copy it}.

\begin{definitionbox}[Pre-Copy Live Migration]
The standard algorithm, used by every platform in \chref{platforms}:

\begin{tight}
  \item \textbf{Round 0.} Copy all of the guest's memory to the destination
        while it keeps running. Track which pages are written during the copy
        --- these are \term{dirty pages}.
  \item \textbf{Round $n$.} Copy the pages dirtied during round $n-1$. The
        guest is still running, so this dirties more.
  \item \textbf{Repeat} until the remaining dirty set is small enough to copy
        within the acceptable pause.
  \item \textbf{Switchover.} Pause the guest, copy the last few pages and the
        processor state, and resume it on the destination. Typically tens of
        milliseconds.
\end{tight}

\smallskip
Each round copies less than the one before, \emph{provided} the guest dirties
memory more slowly than the network transfers it. That proviso is the whole
engineering problem.
\end{definitionbox}

\dsfig{%
\begin{tikzpicture}[font=\scriptsize,
  rb/.style={draw=secondary, fill=secondary!25, line width=0.5pt},
  sw/.style={draw=accent, fill=accent!35, line width=0.7pt},
  lb/.style={font=\scriptsize\bfseries, text=neutral!40!black, anchor=east},
  note/.style={font=\scriptsize\itshape, align=left, anchor=west}
]
  % converging case
  \node[font=\scriptsize\bfseries, text=greenhl!45!black, anchor=west]
       at (0,3.35) {Converges: dirty rate 200 MB/s, link 1 GB/s};
  \node[lb] at (-0.2,2.75) {round 0};
  \draw[rb] (0,2.6) rectangle (8.0,2.92);
  \node[font=\tiny, text=secondary!80!black] at (4.0,2.76) {16 GB, 16.0 s};
  \node[lb] at (-0.2,2.3) {round 1};
  \draw[rb] (0,2.15) rectangle (1.6,2.47);
  \node[font=\tiny, text=secondary!80!black, anchor=west] at (1.75,2.31)
       {3.2 GB, 3.2 s};
  \node[lb] at (-0.2,1.85) {round 2};
  \draw[rb] (0,1.7) rectangle (0.32,2.02);
  \node[font=\tiny, text=secondary!80!black, anchor=west] at (0.47,1.86)
       {640 MB, 0.64 s};
  \node[lb] at (-0.2,1.4) {round 3};
  \draw[rb] (0,1.25) rectangle (0.064,1.57);
  \node[font=\tiny, text=secondary!80!black, anchor=west] at (0.21,1.41)
       {128 MB, 0.13 s};
  \node[lb] at (-0.2,0.95) {switchover};
  \draw[sw] (0,0.8) rectangle (0.08,1.12);
  \node[font=\tiny, text=accent!70!black, anchor=west] at (0.23,0.96)
       {\textbf{paused $\approx$ 50 ms} --- users notice nothing};

  % non-converging case
  \node[font=\scriptsize\bfseries, text=accent!70!black, anchor=west]
       at (0,0.25) {Does not converge: dirty rate 1.2 GB/s, link 1 GB/s};
  \foreach \i/\y in {0/-0.35, 1/-0.8, 2/-1.25, 3/-1.7}{
    \node[lb] at (-0.2,\y+0.16) {round \i};
    \pgfmathsetmacro{\w}{8.0*pow(1.2,\i)}
    \draw[draw=accent, fill=accent!18, line width=0.5pt]
          (0,\y) rectangle (\w,\y+0.32);}
  \node[note, text=accent!70!black] at (10.2,-1.0)
       {each round has \emph{more}\\
        to copy than the last.\\
        It never ends.};

  \node[rounded corners=3pt, draw=primary!50, fill=primary!5, text=primary,
        font=\scriptsize, align=center, inner sep=4pt, text width=12.4cm]
       at (6.0,-2.6)
       {The condition is simply \textbf{dirty rate $<$ transfer rate}. When it
        holds, each round shrinks by the ratio of the two and convergence is
        geometric. When it does not, no amount of patience helps.};
\end{tikzpicture}}%
{Pre-copy migration, drawn to scale for the worked example. The upper case
finishes in 20 seconds with a 50 ms pause; the lower one never finishes at
all.}{pre-copy}

\begin{worked}[Will This Guest Migrate?]
A 16 GB guest is to be migrated over a 10 Gb/s link. Measured: the link
sustains about 1 GB/s after overheads, and the guest dirties memory at 200 MB/s.

\smallskip
\textbf{Step 1 --- the convergence condition.}
\[ \text{dirty rate } 200 \text{ MB/s} \;<\; \text{transfer rate }
   1{,}000 \text{ MB/s} \quad \checkmark \]
Each round will be $200/1000 = 0.2$ of the one before.

\smallskip
\textbf{Step 2 --- round by round.}
\[ \text{round 0: } \frac{16{,}000}{1{,}000} = 16.0 \text{ s},
   \text{ dirtying } 16.0 \times 200 = 3{,}200 \text{ MB} \]
\[ \text{round 1: } 3.2 \text{ s, dirtying } 640 \text{ MB} \qquad
   \text{round 2: } 0.64 \text{ s, dirtying } 128 \text{ MB} \]
\[ \text{round 3: } 0.13 \text{ s, dirtying } 26 \text{ MB} \qquad
   \text{round 4: } 0.026 \text{ s, dirtying } 5 \text{ MB} \]

\smallskip
\textbf{Step 3 --- when to stop.} With a 50 ms pause budget the switchover can
move $0.05 \times 1{,}000 = 50$ MB. Round 4 leaves 5 MB, comfortably under it,
so the hypervisor stops after round 4.

\smallskip
\textbf{Step 4 --- totals.}
\[ \text{elapsed} = 16.0 + 3.2 + 0.64 + 0.13 + 0.026 \approx \mathbf{20
   \text{ s}} \]
\[ \text{transferred} = 16{,}000 + 3{,}200 + 640 + 128 + 26 \approx 20
   \text{ GB} \]
\[ \text{guest paused} = \mathbf{\approx 50 \text{ ms}} \]
A TCP connection does not notice 50 ms. Neither does a user.

\smallskip
\textbf{Step 5 --- now a busy database on the same link.} Same 16 GB, but
dirtying at 1.2 GB/s:
\[ 1{,}200 > 1{,}000 \quad \times \]
Round 0 takes 16 s and dirties 19.2 GB --- \emph{more than the guest has}.
Round 1 is larger than round 0. The migration runs for ever, consuming the link
the whole time.

\smallskip
\textbf{Step 6 --- three ways out.}
\begin{tight}
  \item \textbf{A faster link.} At 25 Gb/s the transfer rate is about 2.5 GB/s
        and $1.2 < 2.5$ converges. This is why migration networks are
        separate and fast.
  \item \textbf{Auto-convergence.} The hypervisor deliberately throttles the
        guest's vCPUs --- 10\%, then 20\%, and so on --- which reduces the
        dirty rate until the condition holds. The guest is slowed on purpose
        to let it move.
  \item \textbf{Post-copy.} Switch over \emph{first}, then fetch pages on
        demand. See below.
\end{tight}

\smallskip
\textbf{Step 7 --- what to take away.} Migration time is governed by memory
size; migration \emph{success} is governed by the ratio of dirty rate to link
speed. A guest that will not migrate is nearly always a write-heavy one on a
link that is too slow, and the fix is usually the network, not the guest.
\end{worked}

\begin{conceptbox}[Post-Copy: the Opposite Risk]
Post-copy inverts the order. The guest is paused almost immediately, its
processor state is sent, and it \emph{resumes on the destination} with almost
no memory. Pages are fetched from the source as the guest faults on them, over
the network.

\smallskip
\textbf{For:} total migration time is bounded and predictable --- every page is
sent exactly once, so there is no convergence problem at all. A guest that
cannot migrate any other way will migrate this way.

\smallskip
\textbf{Against, and it is serious:} between switchover and completion, the
guest's memory lives on \emph{two} hosts. If the network or the source fails in
that window, the guest is unrecoverable --- its pages exist nowhere complete.
Pre-copy's failure mode is ``the migration did not happen''; post-copy's is
``the guest is gone''.

\smallskip
The usual arrangement is hybrid: begin with pre-copy, and if it has not
converged after a few rounds, switch to post-copy for the remainder. That gets
pre-copy's safety for most of the transfer and post-copy's termination
guarantee for the tail.
\end{conceptbox}

\section{High Availability, Fault Tolerance and Quorum}

Live migration is \emph{planned} movement. The unplanned case is a host
failing, and the answer is different.

\term{High availability} restarts the failed host's guests on surviving hosts.
The guests crash and reboot: everything in memory is lost, and the service is
down for as long as a boot takes --- a minute or two. It is cheap, it is
automatic, and it is what almost every cluster runs.

\term{Fault tolerance} runs a shadow copy of the guest on a second host in
lockstep, so a host failure is invisible --- zero downtime, zero data loss. It
costs double the resources, limits the guest's size and vCPU count severely,
and is reserved for the few workloads that genuinely justify it.

\begin{alertbox}[Split brain, and why fencing is not optional]
A cluster decides a host has failed because it stopped answering. But
``stopped answering'' and ``stopped running'' are different things: a network
partition leaves the host alive, still writing to shared storage, and
convinced the others failed.

\smallskip
If the surviving side restarts those guests, the same guest now runs twice,
both instances writing to the same disks. The result is not downtime; it is
\emph{corruption}, and it is worse than the outage that was being avoided.

\smallskip
Two mechanisms prevent it. \term{Quorum}: only a partition holding more than
half the votes may act, so a minority partition stops. This is why clusters
have odd numbers of members or a witness. And \term{fencing}, bluntly named
\term{STONITH} --- shoot the other node in the head: before restarting
anything, the survivors forcibly power off the suspect host through its
management controller or its storage fabric.

\smallskip
A cluster without fencing is not a highly available cluster. It is a system
with one additional, correlated way to lose your data.
\end{alertbox}

\section{RPO, RTO and the Cost of a Nine}

Disaster recovery is about the site, not the host, and two numbers specify it.

\begin{definitionbox}[RPO and RTO]
\term{RPO} --- recovery point objective --- is how much \emph{data} you may
lose, expressed as time. An RPO of one hour means losing up to an hour of work
is acceptable, which means replicating or backing up at least hourly.

\smallskip
\term{RTO} --- recovery time objective --- is how long you may be \emph{down}.
An RTO of four hours means service must be restored within four hours of the
declaration.

\smallskip
They are independent and they cost differently. Nightly backups to tape give an
RPO of 24 hours and an RTO of perhaps a day. Synchronous replication to a
second site gives an RPO near zero --- and an RTO that still depends entirely on
how quickly somebody can bring the second site up, which is a rehearsal
question, not a technology one.
\end{definitionbox}

\begin{worked}[What Each Nine Costs]
A service is offered at ``four nines''. Work out what that means and what the
next nine would require.

\smallskip
\textbf{Step 1 --- the downtime budget.} A year is
$365 \times 24 \times 60 = 525{,}600$ minutes.
\[ 99.9\% : \; 525{,}600 \times 0.001 = 526 \text{ min}
   = \mathbf{8 \text{ h } 46 \text{ min a year}} \]
\[ 99.99\% : \; 52.6 \text{ min a year} = \mathbf{4.4 \text{ min a month}} \]
\[ 99.999\% : \; 5.3 \text{ min a year} = \mathbf{26 \text{ seconds a month}} \]

\smallskip
\textbf{Step 2 --- what fits in each budget.}
\begin{tight}
  \item \textbf{99.9\%}: 8 h 46 min allows a monthly patching window of 40
        minutes and still leaves room for one unplanned incident. Achievable
        with HA and careful scheduling.
  \item \textbf{99.99\%}: 52 minutes a year. A single HA restart of a guest
        costs 1--2 minutes, so you may have about thirty of them --- \emph{and
        no planned downtime at all}. Every patch must be applied by live
        migration.
  \item \textbf{99.999\%}: 5 minutes a year. One HA restart is 20--40\% of the
        annual budget. This requires fault tolerance or clustered application
        software, and it requires that no human decision is in the recovery
        path, because nobody reads an alert in 26 seconds.
\end{tight}

\smallskip
\textbf{Step 3 --- the cost shape.} Each nine divides the budget by ten and
roughly multiplies the cost by three to five: a second site, a second set of
licences, tested automation, rehearsals, and staff who can be woken.

\smallskip
\textbf{Step 4 --- the question to ask instead.} Before accepting an
availability target, ask what the \emph{outage} actually costs per hour. A
service losing \$200 an hour cannot justify the \$150{,}000 a year that five
nines would cost; one losing \$200{,}000 an hour can justify a great deal.
Availability targets chosen without that number are chosen by whoever is most
anxious.

\smallskip
\textbf{Step 5 --- and the measurement everyone forgets.} An availability
figure is meaningless without a stated measurement window and a definition of
``down''. 99.99\% measured monthly permits 4.4 minutes every month; measured
annually it permits 52 minutes in one incident. Those are very different
promises, and the difference is where disputes happen.
\end{worked}

\begin{conceptbox}[A Snapshot Is Not a Backup, Restated]
\chref{storage} made this point about storage; it is worth making again about
disaster recovery, because the failure is so common.

\smallskip
The \term{3-2-1 rule}: three copies of the data, on two different media, with
one off site. A snapshot fails all three --- it is on the same storage, the
same medium, the same site, and it depends on the original.

\smallskip
Ransomware has added a fourth requirement: at least one copy must be
\term{immutable} --- write-once, with a retention lock the backup
administrator's own credentials cannot override. Attackers who gain
administrative access delete the backups first, and an estate whose backup
server trusts the same directory as everything else has one credential between
it and total loss.

\smallskip
And the rule that outranks all of them: \textbf{a backup you have not restored
is a hypothesis}. Test restores on a schedule, and time them --- that
measurement is your real RTO, as opposed to the one in the document.
\end{conceptbox}

\begin{pitfall}
\begin{tight}
  \item \textbf{Expecting any guest to live-migrate.} A write-heavy guest on a
        slow link will not converge. Check the dirty rate against the link
        before promising a maintenance window.
  \item \textbf{Running a cluster without fencing.} Split brain corrupts data,
        which is worse than the outage HA was bought to prevent.
  \item \textbf{Using an even number of cluster members with no witness.} A
        clean split leaves neither side with quorum and both stop.
  \item \textbf{Confusing HA with fault tolerance.} HA restarts the guest; the
        service \emph{is} down and everything in memory is lost.
  \item \textbf{Quoting an availability figure without a window and a
        definition of down.} 99.99\% monthly and 99.99\% annually are different
        promises.
  \item \textbf{Treating replication as backup.} Replication faithfully copies
        the deletion and the encryption to the second site, usually within
        seconds.
  \item \textbf{Never testing a restore.} Then the RTO is a guess, and it is
        always optimistic.
  \item \textbf{Forgetting that a passthrough device prevents migration}
        (\chref{paravirt}). Plan the exception before you need to patch the
        host.
\end{tight}
\end{pitfall}

\begin{lab}[Migrating a Running Guest, and Failing To]
Live-migrates a guest between two hosts with no dropped pings, then makes the
guest dirty memory fast enough that it will not converge. Needs two hosts, or
two nested ones. Expect forty-five minutes.

\begin{lstlisting}[language=bash]
#!/usr/bin/env bash
# ch17_migrate.sh -- pre-copy convergence, observed and then broken.
set -euo pipefail
DEST=${1:?usage: ch17_migrate.sh <destination-host>}

# --- 1. Both hosts must see the same disk. ----------------------------
# Shared storage is the precondition for everything in this chapter.
showmount -e "$DEST" 2>/dev/null || \
  echo "set up NFS or iSCSI first -- see Appendix A"
sudo virsh pool-list --all

# --- 2. Start the guest and give it something to serve. --------------
sudo virsh start vss1 || true ; sleep 25
IP=$(sudo virsh domifaddr vss1 | awk '/ipv4/{split($4,a,"/");print a[1]}')
G () { ssh -o StrictHostKeyChecking=no "ubuntu@$IP" "$1"; }
G "sudo apt-get -qq install -y stress-ng python3 && \
   nohup python3 -m http.server 8000 >/dev/null 2>&1 &"

# --- 3. Watch continuity: ping and HTTP, both running throughout. ----
ping -i 0.2 "$IP" > /tmp/ping.log &
PINGPID=$!
( while true; do curl -s -o /dev/null -w '%{http_code} ' \
    "http://$IP:8000/" ; sleep 0.2; done ) > /tmp/http.log &
CURLPID=$!

# --- 4. Migrate it, and time the whole thing. ------------------------
time sudo virsh migrate --live --verbose --persistent --undefinesource \
  vss1 "qemu+ssh://$DEST/system"

sleep 2 ; kill $PINGPID $CURLPID 2>/dev/null || true
echo "--- dropped pings ---"
grep -c 'icmp_seq' /tmp/ping.log
tail -2 /tmp/ping.log           # look at the loss percentage
echo "--- non-200 responses ---"
tr ' ' '\n' < /tmp/http.log | grep -vc '^200$' || echo 0
# A correct migration drops nothing. The 50 ms pause is inside one
# ping interval.

# --- 5. Now make it refuse to converge. ------------------------------
# stress-ng --vm writes continuously, so the dirty rate exceeds the link.
sudo virsh migrate --live vss1 "qemu+ssh://$DEST/system" || true   # back
sleep 10
G "nohup stress-ng --vm 2 --vm-bytes 75% --vm-keep --timeout 300s \
   >/dev/null 2>&1 &"
sleep 5
echo "--- migrating a guest that dirties faster than the link ---"
sudo virsh migrate --live --verbose --timeout 60 \
  --timeout-suspend vss1 "qemu+ssh://$DEST/system" || \
  echo "did not converge within 60 s -- exactly section 17.1"

# --- 6. Make it converge by throttling the guest on purpose. ---------
sudo virsh migrate --live --verbose --auto-converge \
  --auto-converge-initial 20 --auto-converge-increment 10 \
  vss1 "qemu+ssh://$DEST/system"
# The guest was deliberately slowed so that it could move.

# --- 7. Measure the dirty rate directly, for the arithmetic. ---------
sudo virsh domdirtyrate-calc vss1 --seconds 10 || true
sleep 12 ; sudo virsh domstats vss1 --dirtyrate | grep -i dirty

# --- 8. Unplanned failure: what HA does instead. ---------------------
# No graceful anything: cut the guest off and see what survives.
sudo virsh destroy vss1
echo "the guest is gone; HA would now restart it elsewhere, and"
echo "everything in its memory is lost. That is the difference."

# --- 9. Tear down. ---------------------------------------------------
sudo virsh start vss1 || true
\end{lstlisting}

\textbf{What to notice.} Step 4 should drop no pings at all: the guest moved
between physical machines while serving HTTP and nothing noticed. Step 5 is the
same command failing for the reason the worked example predicts, and step 6 is
the hypervisor buying convergence with the guest's performance.
\end{lab}

\begin{checkpoint}
\begin{tightnum}
  \item A 32 GB guest dirties memory at 400 MB/s and the migration link
        sustains 1.2 GB/s. Give the ratio, the size of the first three rounds,
        and the approximate total elapsed time.
  \item A two-host cluster loses the network between the hosts. Both are alive
        and both can reach the storage. Explain what each host believes, what
        happens without fencing, and what quorum would do differently.
  \item A service promises 99.99\% availability measured annually. Give the
        downtime budget, and say whether a 90-minute planned upgrade once a
        year is compatible with it.
\end{tightnum}
\end{checkpoint}

\begin{chaptersummary}
\begin{tight}
  \item Pre-copy migration copies memory in rounds while the guest runs, then
        pauses briefly to send the remainder. It converges only if the dirty
        rate is below the transfer rate.
  \item A 16 GB guest at 200 MB/s over 1 GB/s converges in about 20 seconds
        with a 50 ms pause. At 1.2 GB/s it never converges at all.
  \item The fixes are a faster link, auto-convergence (throttling the guest on
        purpose), or post-copy --- which always terminates but loses the guest
        if the network fails mid-transfer.
  \item High availability restarts guests after a host failure; fault tolerance
        runs a lockstep shadow. HA loses memory and takes a minute; FT costs
        double.
  \item Quorum and fencing prevent split brain. A cluster without fencing has
        added a way to corrupt data.
  \item RPO is data lost, RTO is time down, and they cost differently. Each
        additional nine divides the budget by ten and multiplies the cost by
        three to five.
  \item 3-2-1, plus one immutable copy, plus tested restores. A backup you have
        not restored is a hypothesis.
\end{tight}
\end{chaptersummary}

\begin{reviewq}
  \item \textbf{(MCQ)} Pre-copy live migration converges when: (a)~the guest
        has fewer than 8 vCPUs; (b)~the dirty rate is lower than the transfer
        rate; (c)~shared storage is used; (d)~the guest is idle.
  \item \textbf{(MCQ)} The characteristic risk of post-copy migration is:
        (a)~it never terminates; (b)~the guest's memory is split across two
        hosts, so a failure in that window loses it; (c)~it needs more
        bandwidth; (d)~it cannot be used with shared storage.
  \item \textbf{(Short)} Explain auto-convergence and what it trades.
  \item \textbf{(Short)} Define RPO and RTO, and give a configuration that
        produces a low RPO and a high RTO.
  \item \textbf{(Short)} State the 3-2-1 rule and explain which parts of it a
        snapshot fails.
  \item \textbf{(Applied)} A 64 GB database guest must be moved for host
        maintenance. It dirties memory at 2.1 GB/s and the migration network
        is 10 Gb/s. Show that it will not converge, give three options with
        their consequences for the database's users, and recommend one.
  \item \textbf{(Applied)} A department is asked to commit to 99.99\%
        availability for a student-facing system. Write the response: the
        downtime budget, what it rules out operationally, what it would cost to
        deliver, and the question you would ask the requester first.
\end{reviewq}
